\documentclass[11pt]{article}
\usepackage[margin=0.9in]{geometry}
\usepackage[dvipsnames]{xcolor}
\usepackage{amsmath}
\usepackage{tikz}
\usepackage{forest}
\usepackage{listings}
\usepackage{needspace}
\setlength{\parindent}{0pt}
\setlength{\parskip}{0.6em}

% ---------- coloured node names inside the text ----------
\newcommand{\rd}[1]{\textcolor{Red}{\textbf{#1}}}   % red node
\newcommand{\bk}[1]{\textbf{#1}}                    % black node

% ---------- tree style ----------
% [5,blackn]  = black node      [20,redn] = red node
% onpath      = thick line from the parent (the path we walked)
% [,phantom]  = empty spot (keeps left/right children in the right place)
\forestset{
	default preamble={
		for tree={circle, draw, minimum size=6.5mm, inner sep=1pt, s sep=4mm, l sep=7mm, font=\footnotesize},
	},
	blackn/.style={fill=black, text=white},
	redn/.style={fill=Red, text=white},
	onpath/.style={edge={very thick}},
}

% ---------- code style ----------
\lstset{
	language=Java,
	basicstyle=\ttfamily\footnotesize,
	keywordstyle=\color{NavyBlue}\bfseries,
	commentstyle=\color{OliveGreen}\itshape,
	stringstyle=\color{Plum},
	showstringspaces=false,
	breaklines=true,
	tabsize=4,
	frame=single,
	columns=fullflexible,
	keepspaces=true
}

\begin{document}
	
	\begin{center}
		{\Large \textbf{Notes on Trees: Binary Search Trees and Red-Black Trees}}
	\end{center}
	
	%%%%%%%%%%%%%%%%%%%%%%%%%%%%%%%%%%%%%%%%%%%%%%%%%%%%%%%%%%%%%%%%%%%%%
	\section{What the lecture said about the test}
	\begin{itemize}
		\item The prof will upload the binary search tree \textbf{code} to GitHub for the test.
		\item The quiz does not ask to display the whole thing, just methods. For example, the code uses \texttt{Stack<>} (in \texttt{displayTree}), and you do not need to know this.
		\item Delete is quite complex, so it is not expected.
		\item Quick sort and partitioning are on Tuesday only.
		\item In any case, know how to code the methods, and make passes through the tree if needed.
	\end{itemize}
	
	%%%%%%%%%%%%%%%%%%%%%%%%%%%%%%%%%%%%%%%%%%%%%%%%%%%%%%%%%%%%%%%%%%%%%
	\section{Trees: the basics}
	\begin{itemize}
		\item A \textbf{graph} stores data in nodes. A node can be anything. Graphs go in many directions: a node can link to many other nodes (or words), and each of those can have plenty of links of its own. Graphs are used a lot in AI, for example in large language models (LLMs).
		\item A \textbf{tree} is a special kind of graph. A tree has one node with no parent, only children. That top node is the \textbf{root}.
		\item A tree where every node has no more than 2 children (0, 1 or 2) is a \textbf{binary tree}.
		\item We pick a key for each node. If the left child has a smaller key and the right child has a key that is greater or equal, it is a \textbf{binary search tree} (BST). The code puts equal keys on the right.
		\item Search takes $\log n$, insertion takes $\log n$, deletion takes $\log n$. This is only possible if the tree is balanced. At worst the tree is just a linked list.
		\item No set size is needed, and you can add nodes as you like. (Lecture note said ``amortization(?)'', not sure what was meant.)
		\item Adding or removing a node can unbalance the tree.
	\end{itemize}
	
	\subsection*{Words to know}
	\begin{itemize}
		\item \textbf{Terminal node} (also called a \textbf{leaf}): a node with no children.
		\item \textbf{Path}: going from the root node to the node we want to reach.
		\item \textbf{Subtree} and \textbf{local root}: a subtree is any node together with everything below it. The top node of that subtree is its local root. (In the code, \texttt{localRoot} is the input of \texttt{preOrder}, \texttt{inOrder} and \texttt{postOrder}.)
	\end{itemize}
	
	\subsection*{Balanced and not balanced}
	A tree that only has left children (or only right children) is biased, so it is not balanced. Its time is $n$ instead of $\log n$.
	
	A tree is \textbf{balanced} when, at every node, the left and right sides differ by no more than one level (each side can have its own levels). If even one node breaks this rule, the whole tree is not balanced.
	% CHANGED: your notes said "no more than one node that is unbalanced", which did not match your other note. I used the one-level rule from your brief notes.
	
	Where $\log n$ comes from: a full tree with $L$ levels holds $2^L-1$ nodes, so $n \approx 2^L$, and
	\[ L = \log_2(n). \]
	So walking from the root down to any node takes about $\log_2(n)$ steps. Below, a balanced tree with 4 levels ($2^4-1 = 15$ nodes):
	
	\begin{center}
		\begin{forest}
			[100
			[50
			[25
			[10]
			[40]
			]
			[75
			[65]
			[80]
			]
			]
			[150
			[125
			[110]
			[130]
			]
			[200
			[175]
			[225]
			]
			]
			]
		\end{forest}
	\end{center}
	% CHANGED: your notes had 90 under 25. 90 is bigger than 50, so it cannot be in the left side of 100's tree. I used 40. Change it if you meant another number.
	
	%%%%%%%%%%%%%%%%%%%%%%%%%%%%%%%%%%%%%%%%%%%%%%%%%%%%%%%%%%%%%%%%%%%%%
	\section{Binary search tree methods, step by step}
	
	\subsection{Insert, minimum and maximum (my recorded code: \texttt{Demo.java})}
	\texttt{insert} starts at the root. If the new key is smaller it goes left, otherwise it goes right, until it reaches an empty spot. My \texttt{Demo} inserts 100, 50, 25, 125, 200, 10:
	\begin{enumerate}
		\item 100: the tree is empty, so 100 is the root.
		\item 50: $50<100$, go left. Empty, place it.
		\item 25: $25<100$ go left, $25<50$ go left. Empty, place it.
		\item 125: $125\ge 100$, go right. Empty, place it.
		\item 200: $200\ge 100$ go right, $200\ge 125$ go right. Empty, place it.
		\item 10: $10<100$ go left, $10<50$ go left, $10<25$ go left. Empty, place it.
	\end{enumerate}
	\Needspace{9\baselineskip}
	The thick line is the path walked for the last insert (10):
	\begin{center}
		\begin{forest}
			[100
			[50, onpath
			[25, onpath
			[10, onpath]
			[,phantom]
			]
			[,phantom]
			]
			[125
			[,phantom]
			[200]
			]
			]
		\end{forest}
	\end{center}
	\begin{itemize}
		\item \texttt{minimum()}: keep going left until there is no left child. The last node reached is the answer: $100 \to 50 \to 25 \to 10$, so the minimum is 10.
		\item \texttt{maximum()}: keep going right until there is no right child: $100 \to 125 \to 200$, so the maximum is 200.
		\item Output of \texttt{Demo}: \texttt{(200, 200.1)(10, 10.1)}. (\texttt{displayNode} in my version uses \texttt{print}, so there is no new line between them.)
	\end{itemize}
	
	\subsection{Find (prof's code: \texttt{TreeApp.java} and \texttt{Tree.java})}
	The prof's \texttt{TreeApp} inserts 50, 25, 75, 12, 37, 43, 30, 33, 87, 93, 97, which gives this tree:
	\begin{center}
		\begin{forest}
			[50
			[25
			[12]
			[37
			[30
			[,phantom]
			[33]
			]
			[43]
			]
			]
			[75
			[,phantom]
			[87
			[,phantom]
			[93
			[,phantom]
			[97]
			]
			]
			]
			]
		\end{forest}
	\end{center}
	Notice this tree is not balanced: at 75 the left side has 0 levels and the right side has 3 levels (75 $\to$ 87 $\to$ 93 $\to$ 97 is almost a linked list).
	
	\Needspace{10\baselineskip}
	\textbf{\texttt{find(33)}}: the thick line is the path.
	\begin{enumerate}
		\item At 50: $33<50$, go left.
		\item At 25: $33>25$, go right.
		\item At 37: $33<37$, go left.
		\item At 30: $33>30$, go right.
		\item At 33: it matches, so \texttt{find} returns this node.
	\end{enumerate}
	If \texttt{current} ever becomes \texttt{null}, the key is not in the tree and \texttt{find} returns \texttt{null}. (\texttt{find} and \texttt{delete} assume the tree is not empty.)
	\begin{center}
		\begin{forest}
			[50
			[25, onpath
			[12]
			[37, onpath
			[30, onpath
			[,phantom]
			[33, onpath]
			]
			[43]
			]
			]
			[75
			[,phantom]
			[87
			[,phantom]
			[93
			[,phantom]
			[97]
			]
			]
			]
			]
		\end{forest}
	\end{center}
	
	\subsection{Delete (not expected on the test, but here are my notes)}
	The prof's \texttt{delete} first finds the node (remembering its parent and whether it is a left child), then there are 3 cases.
	
	\Needspace{9\baselineskip}
	\textbf{Case 1: deleting a node with no children.} The parent just stops pointing to it. (If it was the root, the tree becomes empty.) Delete 71:
	\begin{center}
		\begin{forest}
			[80
			[52
			[48]
			[71, onpath]
			]
			[,phantom]
			]
		\end{forest}
		\quad$\Rightarrow$\quad
		\begin{forest}
			[80
			[52
			[48]
			[,phantom]
			]
			[,phantom]
			]
		\end{forest}
	\end{center}
	
	\Needspace{9\baselineskip}
	\textbf{Case 2: deleting a node with one child.} The child (with everything below it) takes its place. In the code: if there is no right child, replace with the left subtree, and if there is no left child, replace with the right subtree. Delete 80 (the root), which only has the left child 52:
	\begin{center}
		\begin{forest}
			[80
			[52, onpath
			[48]
			[71]
			]
			[,phantom]
			]
		\end{forest}
		\quad$\Rightarrow$\quad
		\begin{forest}
			[52
			[48]
			[71]
			]
		\end{forest}
	\end{center}
	
	\textbf{Case 3: deleting a node with two children} (my notes said ``three children'', but it is two).
	% CHANGED: "three children" -> "two children" (a binary tree node has at most two).
	Find the \textbf{successor}: the next-highest key. Go to the right child, then keep going left while there is a left child. The last node reached is the smallest value in the right subtree.
	
	\emph{From the lecture:} ``Go right from 175 to find 180. See if 180 has children of its own. If it doesn't, then it is the successor, it replaces the 175 with 180 and returns 175.''
	
	\Needspace{12\baselineskip}
	\textbf{Worked example:} delete 50 from the 100-tree (the parent of 50 is 100). Thick line = the walk to the successor.
	\begin{center}
		\begin{forest}
			[50
			[25
			[10]
			[40]
			]
			[75, onpath
			[65, onpath]
			[80]
			]
			]
		\end{forest}
		\quad$\Rightarrow$\quad
		\begin{forest}
			[65
			[25
			[10]
			[40]
			]
			[75
			[,phantom]
			[80]
			]
			]
		\end{forest}
	\end{center}
	\begin{enumerate}
		\item Find 50: at 100, $50<100$, go left. Found. Its parent is 100, and it is a left child.
		\item 50 has two children, so call \texttt{getSuccessor}. Go to the right child, 75.
		\item Keep going left: 75 has a left child, 65. 65 has no left child, so stop. The successor is 65 and its parent is 75.
		\item 65 is not 50's right child, so fix the links: 75's left link becomes 65's right child (empty), and 65's right link becomes 75.
		\item Connect 100's left link to 65 instead of 50, and give 65 the left child of 50 (25).
	\end{enumerate}
	
	%%%%%%%%%%%%%%%%%%%%%%%%%%%%%%%%%%%%%%%%%%%%%%%%%%%%%%%%%%%%%%%%%%%%%
	\section{Traversals}
	A traversal visits every node once. In the code, each one calls itself on a subtree. \texttt{localRoot} is the node we are at, and if it is \texttt{null} we do nothing and go back.
	
	\begin{itemize}
		\item \textbf{Preorder} (\texttt{traverse(1)}):
		\begin{enumerate}
			\item Visit the node (print it).
			\item Call itself to traverse the node's left subtree.
			\item Call itself to traverse the node's right subtree.
		\end{enumerate}
		\item \textbf{In-order} (\texttt{traverse(2)}): left subtree, then visit the node, then right subtree. Start at the bottom left and go up. For 3 nodes (a root with two children), the left child is 1, the root is 2 and the right child is 3.
		\item \textbf{Postorder} (\texttt{traverse(3)}): left subtree, then right subtree, then visit the node last.
		% CHANGED: your notes said postorder does "the same in the opposite order". The prof's postOrder code is left, right, visit.
	\end{itemize}
	
	\Needspace{7\baselineskip}
	Results on the prof's tree (the same tree as in the find example):
	\begin{center}
		\begin{tabular}{ll}
			Preorder  & 50 25 12 37 30 33 43 75 87 93 97 \\
			In-order  & 12 25 30 33 37 43 50 75 87 93 97 \\
			Postorder & 12 33 30 43 37 25 97 93 87 75 50
		\end{tabular}
	\end{center}
	
	\textbf{In-order walk, step by step:}
	\begin{enumerate}
		\item Start at 50, go left to 25, go left to 12.
		\item 12 has no left child, so print \textbf{12}. It has no right child, so go back up to 25.
		\item Print \textbf{25}. Go right to 37, then left to 30.
		\item 30 has no left child, so print \textbf{30}. Go right to 33 and print \textbf{33}.
		\item Back up to 37 and print \textbf{37}. Go right to 43 and print \textbf{43}.
		\item The left side of 50 is done. Print \textbf{50}.
		\item Go right to 75. It has no left child, so print \textbf{75}. Then 87, 93 and 97 are each printed on the way down, because none of them has a left child.
	\end{enumerate}
	In a binary search tree, in-order gives the keys in sorted order.
	
	\textbf{Example with an expression tree} (in-order, preorder and postorder give infix, prefix and postfix):
	\begin{center}
		\begin{forest}
			[$\cdot$
			[A]
			[$+$
			[B]
			[C]
			]
			]
		\end{forest}
	\end{center}
	% CHANGED: infix was written A (C+C) in your notes; the tree has B and C, so I wrote A (B+C).
	\begin{itemize}
		\item Infix (in-order): $A \cdot (B+C)$
		\item Prefix (preorder): $\cdot\, A + B C$
		\item Postfix (postorder): $A B C + \cdot$
	\end{itemize}
	
	%%%%%%%%%%%%%%%%%%%%%%%%%%%%%%%%%%%%%%%%%%%%%%%%%%%%%%%%%%%%%%%%%%%%%
	\section{Rotations}
	\textbf{Right rotation.} Raise the left child, and drop the node down to the right. Turn this
	\begin{center}
		\begin{forest}
			[50
			[25
			[10]
			[,phantom]
			]
			[,phantom]
			]
		\end{forest}
	\end{center}
	into this
	\begin{center}
		\begin{forest}
			[25
			[10]
			[50]
			]
		\end{forest}
	\end{center}
	Doing it in the other direction is a \textbf{left rotation} (raise the right child, drop the node down to the left). If the raised node already had a child on the side where the dropped node goes, that child moves over to the other side of the dropped node (it ``crosses over'').
	
	\textbf{Double rotation.} If the nodes make a zig-zag instead of a straight line, one rotation is not enough, so we do two. This is done to make sure our nodes stay balanced. (Next lecture: keeping a tree balanced.)
	
	\textbf{Example 1: 50, 25, 30} (25 is left of 50, 30 is right of 25).
	Left rotation at 25 first (raise 30, drop 25) so that everything goes left in a straight line. Then a right rotation at 50 (raise 30). Now 30 is in the middle, 25 on the left and 50 on the right.
	\begin{center}
		\begin{forest}
			[50
			[25
			[,phantom]
			[30]
			]
			[,phantom]
			]
		\end{forest}
		\quad$\Rightarrow$\quad
		\begin{forest}
			[50
			[30
			[25]
			[,phantom]
			]
			[,phantom]
			]
		\end{forest}
		\quad$\Rightarrow$\quad
		\begin{forest}
			[30
			[25]
			[50]
			]
		\end{forest}
	\end{center}
	% CHANGED: your notes said "50 30 20"; I used 50 30 25 (the 25 from the example).
	
	\textbf{Example 2: 50, 75, 60} (75 is right of 50, 60 is left of 75).
	Right rotation at 75 first (raise 60, drop 75) so that everything goes right in a straight line. Then a left rotation at 50 (raise 60). Now 60 is on top, with children 50 and 75.
	\begin{center}
		\begin{forest}
			[50
			[,phantom]
			[75
			[60]
			[,phantom]
			]
			]
		\end{forest}
		\quad$\Rightarrow$\quad
		\begin{forest}
			[50
			[,phantom]
			[60
			[,phantom]
			[75]
			]
			]
		\end{forest}
		\quad$\Rightarrow$\quad
		\begin{forest}
			[60
			[50]
			[75]
			]
		\end{forest}
	\end{center}
	
	%%%%%%%%%%%%%%%%%%%%%%%%%%%%%%%%%%%%%%%%%%%%%%%%%%%%%%%%%%%%%%%%%%%%%
	\section{Red-black trees}
	
	\subsection{The 4 axioms}
	\begin{enumerate}
		\item Every node has a colour, either red or black.
		\item The root node must be black. If it is red, change it to black.
		\item A red node cannot have a red child.
		\item The black height has to be the same across the whole tree.
	\end{enumerate}
	For axiom 4: at every branch, count the black nodes along the path from the root (counting the root). Each path must have the same number of black nodes. If it is not the same at any branch, it is not a red-black tree.
	
	\textbf{Null children.} A terminal node (leaf) is a node with no children. When we count along a path, we go all the way down to the \emph{null children}, not just to the terminal nodes. A leaf has two null children, a node with one child has one null child, and a node with two children has no null children. Null children are usually counted as black, but it does not matter: every path gets that extra null child, so the black heights stay equal anyway.
	
	\textbf{Why it stays $\log n$.} Imagine a red-black tree with 4 black nodes in a row down the left side. To have the same black height on the right side while fitting as many nodes as possible, the right side alternates between black and red: \bk{black}, \rd{red}, \bk{black}, \rd{red}, and so on. That doubles the number of nodes on that path (8 instead of 4). So in the worst case, one side is no more than twice as long as the other. This is still $\log n$ (even if it is technically $2\log n$).
	
	We can either change the colour, or rotate. We use the colour as an indicator for when to rotate.
	
	\subsection{How to add a node}
	\begin{enumerate}
		\item Walk down from the root like a normal BST insert (smaller $\to$ left, otherwise $\to$ right) until the spot is empty. Place the new node there.
		\item The new node is always \rd{red}.
		\item If it is the root, colour it \bk{black}. If its parent is black, we are done.
		\item If its parent is red, we have a red parent with a red child, which is not allowed. Look at the \textbf{uncle} (the parent's sibling):
		\begin{itemize}
			\item \textbf{Uncle is \rd{red}:} bring the black down from the grandparent and give it to the parent and the uncle (both become black). The grandparent becomes red. Then check the grandparent: if it is the root, change it to black. If its own parent is red, it is a new problem, so repeat from this step.
			\item \textbf{Uncle is \bk{black} (or empty, null counts as black):} we cannot bring the black down, so rotate. This is the only time we rotate.
			\begin{itemize}
				\item Parent and new node on the same side (a straight line): one rotation at the grandparent. Right rotation if both are left children, left rotation if both are right children.
				\item Parent and new node on different sides (a zig-zag): double rotation.
				\item After rotating, the top node is black and the two nodes below it are red. The colour stays with the place, not with the number.
			\end{itemize}
		\end{itemize}
	\end{enumerate}
	
	\subsection{Example: insert 5, 20, 1, 10, 4, 6, 7, 3, 8, 9 (a random order)}
	Red nodes are filled red and black nodes are filled black. The thick line is the path walked from the root down to the new node. In the text, node names are coloured the same way.
	
	%%%%%%%%%%%%%%%%%%%%%%%%%%%%%%%%%%%%%%%%%%%%%%%%%%%%%%%%%%%%%%%%%%%%%
	\Needspace{17\baselineskip}
	\subsubsection*{Step 1: insert 5}
	The tree is empty, so 5 is the root. New nodes start red, but the root must be black, so colour it \bk{black}.
	\begin{center}
		\begin{forest}
			[5, blackn]
		\end{forest}
	\end{center}
	
	%%%%%%%%%%%%%%%%%%%%%%%%%%%%%%%%%%%%%%%%%%%%%%%%%%%%%%%%%%%%%%%%%%%%%
	\Needspace{17\baselineskip}
	\subsubsection*{Step 2: insert 20}
	Walk: at \bk{5}, $20>5$ so go right. Empty, so place 20 as \rd{red}. The parent \bk{5} is black, so there is no problem.
	\begin{center}
		\begin{forest}
			[5, blackn
			[,phantom]
			[20, redn, onpath]
			]
		\end{forest}
	\end{center}
	
	%%%%%%%%%%%%%%%%%%%%%%%%%%%%%%%%%%%%%%%%%%%%%%%%%%%%%%%%%%%%%%%%%%%%%
	\Needspace{17\baselineskip}
	\subsubsection*{Step 3: insert 1}
	Walk: at \bk{5}, $1<5$ so go left. Empty, so place 1 as \rd{red}. The parent \bk{5} is black, so there is no problem.
	\begin{center}
		\begin{forest}
			[5, blackn
			[1, redn, onpath]
			[20, redn]
			]
		\end{forest}
	\end{center}
	% CHANGED: your notes said 20 and 1 were black here. New nodes start red, and step 4 (uncle 1 is red, parent 20 is red) only works if they are red.
	
	%%%%%%%%%%%%%%%%%%%%%%%%%%%%%%%%%%%%%%%%%%%%%%%%%%%%%%%%%%%%%%%%%%%%%
	\Needspace{17\baselineskip}
	\subsubsection*{Step 4: insert 10}
	Walk: at \bk{5}, $10>5$ so go right $\to$ at \rd{20}, $10<20$ so go left $\to$ empty, place 10 as \rd{red}.
	
	Problem: the red parent \rd{20} has a red child \rd{10}. Look at the uncle (the parent's sibling): it is \rd{1}, which is red. Both the parent and the uncle are red, so bring the black down from the grandparent \bk{5} and give it to 20 and 1. Now 20 and 1 are black, 10 stays red, but the root 5 became red. The root must be black, so change it back to black. Solved!
	\begin{center}
		\begin{forest}
			[5, blackn
			[1, redn]
			[20, redn, onpath
			[10, redn, onpath]
			[,phantom]
			]
			]
		\end{forest}
		\quad$\Rightarrow$\quad
		\begin{forest}
			[5, redn
			[1, blackn]
			[20, blackn
			[10, redn]
			[,phantom]
			]
			]
		\end{forest}
		\quad$\Rightarrow$\quad
		\begin{forest}
			[5, blackn
			[1, blackn]
			[20, blackn
			[10, redn]
			[,phantom]
			]
			]
		\end{forest}
	\end{center}
	
	%%%%%%%%%%%%%%%%%%%%%%%%%%%%%%%%%%%%%%%%%%%%%%%%%%%%%%%%%%%%%%%%%%%%%
	\Needspace{17\baselineskip}
	\subsubsection*{Step 5: insert 4}
	Walk: at \bk{5}, $4<5$ so go left $\to$ at \bk{1}, $4>1$ so go right $\to$ empty, place 4 as \rd{red}. The parent \bk{1} is black, so there is no problem (a black node with a red child is fine).
	\begin{center}
		\begin{forest}
			[5, blackn
			[1, blackn, onpath
			[,phantom]
			[4, redn, onpath]
			]
			[20, blackn
			[10, redn]
			[,phantom]
			]
			]
		\end{forest}
	\end{center}
	
	%%%%%%%%%%%%%%%%%%%%%%%%%%%%%%%%%%%%%%%%%%%%%%%%%%%%%%%%%%%%%%%%%%%%%
	\Needspace{17\baselineskip}
	\subsubsection*{Step 6: insert 6}
	Walk: at \bk{5}, $6>5$ so go right $\to$ at \bk{20}, $6<20$ so go left $\to$ at \rd{10}, $6<10$ so go left $\to$ empty, place 6 as \rd{red}.
	
	Problem: red parent \rd{10}, red child \rd{6}. Look at the uncle (the other child of 20): there is none, and an empty spot counts as black. A black uncle means we cannot bring the black down, so we must rotate. The parent (10) and the new node (6) are both left children, so they are in a straight line: one \textbf{right rotation} at the grandparent 20. Raise 10 and drop 20. The colours stay with the places, so the top is black and the two below it are red.
	\begin{center}
		\begin{forest}
			[5, blackn
			[1, blackn
			[,phantom]
			[4, redn]
			]
			[20, blackn, onpath
			[10, redn, onpath
			[6, redn, onpath]
			[,phantom]
			]
			[,phantom]
			]
			]
		\end{forest}
		\quad$\Rightarrow$\quad
		\begin{forest}
			[5, blackn
			[1, blackn
			[,phantom]
			[4, redn]
			]
			[10, blackn
			[6, redn]
			[20, redn]
			]
			]
		\end{forest}
	\end{center}
	
	%%%%%%%%%%%%%%%%%%%%%%%%%%%%%%%%%%%%%%%%%%%%%%%%%%%%%%%%%%%%%%%%%%%%%
	\Needspace{17\baselineskip}
	\subsubsection*{Step 7: insert 7}
	Walk: at \bk{5}, $7>5$ so go right $\to$ at \bk{10}, $7<10$ so go left $\to$ at \rd{6}, $7>6$ so go right $\to$ empty, place 7 as \rd{red}.
	
	Problem: red parent \rd{6}, red child \rd{7}. The uncle (the sibling of 6) is \rd{20}, which is red. Use the uncle rule: the grandparent 10 becomes red, and 6 and 20 become black. The parent of 10 is the black 5, so we are fine. It is still balanced, so it remains $\log n$.
	\begin{center}
		\begin{forest}
			[5, blackn
			[1, blackn
			[,phantom]
			[4, redn]
			]
			[10, blackn, onpath
			[6, redn, onpath
			[,phantom]
			[7, redn, onpath]
			]
			[20, redn]
			]
			]
		\end{forest}
		\quad$\Rightarrow$\quad
		\begin{forest}
			[5, blackn
			[1, blackn
			[,phantom]
			[4, redn]
			]
			[10, redn
			[6, blackn
			[,phantom]
			[7, redn]
			]
			[20, blackn]
			]
			]
		\end{forest}
	\end{center}
	
	%%%%%%%%%%%%%%%%%%%%%%%%%%%%%%%%%%%%%%%%%%%%%%%%%%%%%%%%%%%%%%%%%%%%%
	\Needspace{17\baselineskip}
	\subsubsection*{Step 8: insert 3}
	Walk: at \bk{5}, $3<5$ so go left $\to$ at \bk{1}, $3>1$ so go right $\to$ at \rd{4}, $3<4$ so go left $\to$ empty, place 3 as \rd{red}.
	
	Problem: red parent \rd{4}, red child \rd{3}. The uncle (the left child of 1) is empty, which counts as black, so we rotate. The parent 4 is a right child and the new node 3 is a left child. They are on different sides (a zig-zag), so we need a \textbf{double rotation}:
	\begin{enumerate}
		\item Right rotation at 4: raise 3, drop 4. Now 1, 3, 4 are in a straight line going right.
		\item Left rotation at 1: raise 3, drop 1 (and 4 stays on the right).
	\end{enumerate}
	Then colour: 3 on top is black, 1 and 4 below are red.
	\begin{center}
		\begin{forest}
			[5, blackn
			[1, blackn, onpath
			[,phantom]
			[4, redn, onpath
			[3, redn, onpath]
			[,phantom]
			]
			]
			[10, redn
			[6, blackn
			[,phantom]
			[7, redn]
			]
			[20, blackn]
			]
			]
		\end{forest}
	\end{center}
	\Needspace{9\baselineskip}
	The part being fixed, after each rotation:
	\begin{center}
		\begin{forest}
			[1, blackn
			[,phantom]
			[3, redn
			[,phantom]
			[4, redn]
			]
			]
		\end{forest}
		\quad$\Rightarrow$\quad
		\begin{forest}
			[3, blackn
			[1, redn]
			[4, redn]
			]
		\end{forest}
	\end{center}
	\Needspace{9\baselineskip}
	The whole tree now (still balanced):
	\begin{center}
		\begin{forest}
			[5, blackn
			[3, blackn
			[1, redn]
			[4, redn]
			]
			[10, redn
			[6, blackn
			[,phantom]
			[7, redn]
			]
			[20, blackn]
			]
			]
		\end{forest}
	\end{center}
	
	%%%%%%%%%%%%%%%%%%%%%%%%%%%%%%%%%%%%%%%%%%%%%%%%%%%%%%%%%%%%%%%%%%%%%
	\Needspace{17\baselineskip}
	\subsubsection*{Step 9: insert 8}
	Walk: at \bk{5}, $8>5$ so go right $\to$ at \rd{10}, $8<10$ so go left $\to$ at \bk{6}, $8>6$ so go right $\to$ at \rd{7}, $8>7$ so go right $\to$ empty, place 8 as \rd{red}.
	
	Problem: red parent \rd{7}, red child \rd{8}. The uncle (the left child of 6) is empty, so it is black, so we rotate. The parent 7 and the new node 8 are both right children (a straight line), so one \textbf{left rotation} at the grandparent 6: raise 7, drop 6, and 8 stays with 7. Now 7 is the parent with children 6 and 8, both red. Counting the empty spots too, it is still balanced.
	\begin{center}
		\begin{forest}
			[5, blackn
			[3, blackn
			[1, redn]
			[4, redn]
			]
			[10, redn, onpath
			[6, blackn, onpath
			[,phantom]
			[7, redn, onpath
			[,phantom]
			[8, redn, onpath]
			]
			]
			[20, blackn]
			]
			]
		\end{forest}
		\quad$\Rightarrow$\quad
		\begin{forest}
			[5, blackn
			[3, blackn
			[1, redn]
			[4, redn]
			]
			[10, redn
			[7, blackn
			[6, redn]
			[8, redn]
			]
			[20, blackn]
			]
			]
		\end{forest}
	\end{center}
	
	%%%%%%%%%%%%%%%%%%%%%%%%%%%%%%%%%%%%%%%%%%%%%%%%%%%%%%%%%%%%%%%%%%%%%
	\Needspace{17\baselineskip}
	\subsubsection*{Step 10: insert 9}
	Walk: at \bk{5}, $9>5$ so go right $\to$ at \rd{10}, $9<10$ so go left $\to$ at \bk{7}, $9>7$ so go right $\to$ at \rd{8}, $9>8$ so go right $\to$ empty, place 9 as \rd{red}.
	
	\Needspace{10\baselineskip}
	\textbf{Problem 1:} red parent \rd{8}, red child \rd{9}. The uncle (the sibling of 8) is \rd{6}, which is red. Use the uncle rule: 6 and 8 become black, and the grandparent 7 becomes red.
	\begin{center}
		\begin{forest}
			[5, blackn
			[3, blackn
			[1, redn]
			[4, redn]
			]
			[10, redn, onpath
			[7, blackn, onpath
			[6, redn]
			[8, redn, onpath
			[,phantom]
			[9, redn, onpath]
			]
			]
			[20, blackn]
			]
			]
		\end{forest}
		\quad$\Rightarrow$\quad
		\begin{forest}
			[5, blackn
			[3, blackn
			[1, redn]
			[4, redn]
			]
			[10, redn
			[7, redn
			[6, blackn]
			[8, blackn
			[,phantom]
			[9, redn]
			]
			]
			[20, blackn]
			]
			]
		\end{forest}
	\end{center}
	
	\Needspace{12\baselineskip}
	\textbf{Problem 2:} that created a new one: 7 and 10 are both red. The uncle of 7 (the sibling of 10) is \bk{3}, which is black. The uncle and parent do not agree, so rotate. 10 is a right child and 7 is a left child (a zig-zag), so a \textbf{double rotation}:
	\begin{enumerate}
		\item Right rotation at 10: raise 7, drop 10. The old 8 does not fit on the right of 7 anymore (10 is there now), so it has to move to where it fits: bigger than 7 and smaller than 10, so it becomes the left child of 10 (a crossover).
		\item Left rotation at 5: raise 7, drop 5. Now 6 does not fit, so it moves to where it fits: bigger than 5 and smaller than 7, so it becomes the right child of 5.
	\end{enumerate}
	Then colour: 7 on top is black, 5 and 10 below are red.
	\begin{center}
		\begin{forest}
			[5, blackn
			[3, blackn
			[1, redn]
			[4, redn]
			]
			[7, redn
			[6, blackn]
			[10, redn
			[8, blackn
			[,phantom]
			[9, redn]
			]
			[20, blackn]
			]
			]
			]
		\end{forest}
		\quad$\Rightarrow$\quad
		\begin{forest}
			[7, blackn
			[5, redn
			[3, blackn
			[1, redn]
			[4, redn]
			]
			[6, blackn]
			]
			[10, redn
			[8, blackn
			[,phantom]
			[9, redn]
			]
			[20, blackn]
			]
			]
		\end{forest}
	\end{center}
	The left picture is after the right rotation at 10. The right picture is after the left rotation at 5 and the colouring, which is the final tree.
	
	\textbf{Check the black height.} Every path from the root (7) down has 2 black nodes, counting the root:
	\begin{itemize}
		\item \bk{7} $\to$ \rd{5} $\to$ \bk{3} $\to$ \rd{1}: 2 black nodes
		\item \bk{7} $\to$ \rd{5} $\to$ \bk{3} $\to$ \rd{4}: 2 black nodes
		\item \bk{7} $\to$ \rd{5} $\to$ \bk{6}: 2 black nodes
		\item \bk{7} $\to$ \rd{10} $\to$ \bk{8} $\to$ \rd{9}: 2 black nodes
		\item \bk{7} $\to$ \rd{10} $\to$ \bk{8} (and then its empty left spot): 2 black nodes
		\item \bk{7} $\to$ \rd{10} $\to$ \bk{20}: 2 black nodes
	\end{itemize}
	No red node has a red child, the root is black, and the black height is the same everywhere, so it is a valid red-black tree. Still balanced and equal black height.
	
	%%%%%%%%%%%%%%%%%%%%%%%%%%%%%%%%%%%%%%%%%%%%%%%%%%%%%%%%%%%%%%%%%%%%%
	\subsection{Now imagine a non-random set of numbers: 1, 2, 3, 4, 5, 6, 7, 8, 9, 10}
	(No steps in my notes yet. Probably the point: with numbers in order, a plain binary search tree becomes a chain like a linked list, but a red-black tree keeps rotating so it stays balanced.)
	
	Binary search tree \& red-black tree: the prof will upload the binary search tree code to GitHub for the test.
	
	%%%%%%%%%%%%%%%%%%%%%%%%%%%%%%%%%%%%%%%%%%%%%%%%%%%%%%%%%%%%%%%%%%%%%
	\newpage
	\section{Code}
	
	\subsection{What is different between the two versions}
	\begin{itemize}
		\item \textbf{My recorded lecture code:} \texttt{Demo}, \texttt{Node}, \texttt{Tree} with only \texttt{insert}, \texttt{minimum} and \texttt{maximum}. \texttt{displayNode} prints \texttt{(key, data)} with no new line.
		\item \textbf{Prof's code (uploaded afterwards):} adds \texttt{find}, \texttt{delete} with \texttt{getSuccessor}, \texttt{traverse} (\texttt{preOrder}, \texttt{inOrder}, \texttt{postOrder}) and \texttt{displayTree} (uses \texttt{Stack}, not needed for the quiz). It still has \texttt{minimum} and \texttt{maximum}. \texttt{TreeApp} is a menu to try it all, and \texttt{APIDemo} uses Java's built-in \texttt{TreeMap} and \texttt{TreeSet} (Java builds these on red-black trees).
		\item \texttt{insert} is the same in both versions.
	\end{itemize}
	
	\subsection{My recorded lecture code}
	
	\textbf{Demo.java}
	\begin{lstlisting}
		package p1;
		
		public class Demo {
			public static void main(String[] args){
				Tree myTree= new Tree();
				myTree.insert(100,100.1);
				myTree.insert(50,50.1);
				myTree.insert(25,125.1);
				myTree.insert(125,120.1);
				myTree.insert(200,200.1);
				myTree.insert(10,10.1);
				myTree.maximum().displayNode();
				myTree.minimum().displayNode();
			}
		}
	\end{lstlisting}
	
	\textbf{Node.java}
	\begin{lstlisting}
		package p1;
		
		public class Node {
			public int iData; // data item key
			public double dData;
			public Node leftChild;
			public Node rightChild;
			
			public void displayNode(){
				System.out.print("(" + iData);
				System.out.print(", ");
				System.out.print(dData + ")");
			}
		}
	\end{lstlisting}
	
	\textbf{Tree.java}
	\begin{lstlisting}
		package p1;
		
		public class Tree {
			private Node root; //first node of the tree
			
			public Tree(){
				root = null;
			}
			
			
			public Node minimum(){
				Node current = root;
				Node last = null;
				while (current != null){
					last = current;
					current = current.leftChild;
				}
				return last;
			}
			
			public Node maximum(){
				Node current = root;
				Node last = null;
				while (current!= null){
					last = current;
					current = current.rightChild;
				}
				return last;
			}
			
			public void insert(int id, double dd){
				Node newNode = new Node();
				newNode.iData = id;
				newNode.dData = dd;
				
				if(root == null){
					root = newNode;
				}else{
					Node current = root;
					Node parent;
					while (true){
						parent = current;
						if (id < current.iData){//go left
							current = current.leftChild;
							if (current == null){//reached the end of the line
								parent.leftChild = newNode;
								return;
							}
						}
						
						else{
							current = current.rightChild;
							if (current == null){
								parent.rightChild = newNode;
								return;
								
							}
						}
					}
				}
			}
			
		}
	\end{lstlisting}
	
	\subsection{Prof's code}
	
	\textbf{APIDemo.java}
	
	Trace: \texttt{TreeSet} keeps its keys sorted and ignores a repeated key, so the loop prints A F G K M N S Y (the second \texttt{"A"} is not added again) and \texttt{contains("A")} is \texttt{true}. In the commented-out \texttt{TreeMap}, the second \texttt{put("Z", 100)} replaces the first value, so \texttt{get("Z")} would print 100.
	\begin{lstlisting}
		package p1;
		import java.util.TreeMap;
		import java.util.TreeSet;
		
		public class APIDemo {
			
			public static void main(String[] args) {
				//        TreeMap<String, Integer> theMap = new TreeMap<>();
				//        theMap.put("M", 100);
				//        theMap.put("G", 200);
				//        theMap.put("Z", 50);
				//        theMap.put("K", 300);
				//        theMap.put("B", 600);
				//        theMap.put("Z", 100);
				//        System.out.println(theMap.get("Z"));
				
				
				TreeSet<String> theTree = new TreeSet<>();
				theTree.add("M");
				theTree.add("F");
				theTree.add("G");
				theTree.add("A");
				theTree.add("K");
				theTree.add("S");
				theTree.add("Y");
				theTree.add("N");
				theTree.add("A");
				
				
				for(String s : theTree) {
					System.out.println(s);
				}
				System.out.println(theTree.contains("A"));
			}
			
		}
	\end{lstlisting}
	
	\textbf{Node.java}
	\begin{lstlisting}
		package p1;
		
		public class Node {
			
			public int iData;        // data item (key)
			public double dData;     // data item
			
			public Node leftChild;       // this node's left child
			public Node rightChild;      // this node's right child
			
			public void displayNode()     // display ourselves
			{
				System.out.print(iData);
				System.out.print(", ");
				System.out.println(dData);
			}
		}
	\end{lstlisting}
	
	\textbf{Tree.java}
	\begin{lstlisting}
		package p1;
		import java.util.Stack;
		
		public class Tree {
			private Node root; // first node of tree
			
			public Tree() { // constructor
				root = null; // no nodes in tree yet
			}
			
			public Node find(int key) { // find node with given key
				// (assumes non-empty tree)
				Node current = root; // start at root
				while (current.iData != key) { // while no match;
					
					if (key < current.iData) { // go left ?
						
						current = current.leftChild;
					} else {// or go right?
						
						current = current.rightChild;
					}
					if (current == null) { // if no child, did't find it.
						return null;
					}
				} // end while()
				return current; // found it
			} // end find()
			
			public void insert(int id, double dd) {
				Node newNode = new Node(); // make new node
				newNode.iData = id; // insert data
				newNode.dData = dd;
				
				if (root == null) { // no node in root
					root = newNode;
				} else {
					Node current = root; // start at root
					Node parent;
					while (true) {
						parent = current;
						if (id < current.iData) { // go left?
							current = current.leftChild;
							if (current == null) { // if end of the line
								parent.leftChild = newNode;
								return;
							}
						} // end if go left
						else { // or go right?
							current = current.rightChild;
							if (current == null) { // if end of the line
								parent.rightChild = newNode; // insert on right
								return;
							}
						} // end else go right
					} // end while
				} // end else not root
			} // end insert()
			
			public boolean delete(int key) { // delete node with given key
				// (assumes non-empty list)
				Node current = root;
				Node parent = root;
				boolean isLeftChild = true;
				
				while (current.iData != key) { // search for key
					parent = current;
					if (key < current.iData) { // go left
						isLeftChild = true;
						current = current.leftChild;
					} else { // or go right
						isLeftChild = false;
						current = current.rightChild;
					}
					
					if (current == null) { // end of the line
						return false;
					}
				} // end while
				
				// found node to delete, it is just a matter of deleting it from this
				// point on
				
				// if it has no children, simply delete it
				if (current.leftChild == null && current.rightChild == null) {
					if (current == root) { // if root,
						root = null; // tree is empty
					} else if (isLeftChild) {
						parent.leftChild = null; // disconnect from parent
					} else {
						parent.rightChild = null; // disconnect from parent
					}
				}
				
				// if no right child, replace with left subtree
				else if (current.rightChild == null) {
					if (current == root) {
						root = current.leftChild;
					} else if (isLeftChild) {
						parent.leftChild = current.leftChild;
					} else {
						parent.rightChild = current.leftChild;
					}
				}
				// if no left child, replace with right subtree
				else if (current.leftChild == null) {
					if (current == root) {
						root = current.rightChild;
					} else if (isLeftChild) {
						parent.leftChild = current.rightChild;
					} else {
						parent.rightChild = current.rightChild;
					}
				} else { // two children
					// get successor of node to delete (current)
					Node successor = getSuccessor(current);
					// connect parent of current to successor instead
					if (current == root) {
						root = successor;
					} else if (isLeftChild) {
						parent.leftChild = successor;
					} else {
						parent.rightChild = successor;
					}
					// connect successor to current's left child
					successor.leftChild = current.leftChild;
				} // end else two children
				
				// (successor cannot have a left child)
				return true;
			} // end of delete
			
			// returns node with next-highest value after delNode
			// goes to right child then right child's left descendants
			private Node getSuccessor(Node delNode) {
				Node successorParent = delNode;
				Node successor = delNode;
				Node current = delNode.rightChild; // go to right child
				while (current != null) { // until no more left children
					successorParent = successor;
					successor = current;
					current = current.leftChild; // go to left child
				}
				
				if (successor != delNode.rightChild) { // if successor not right child,
					// make connections
					successorParent.leftChild = successor.rightChild;
					successor.rightChild = delNode.rightChild;
				}
				return successor;
			}
			
			public void traverse(int traverseType) {
				switch (traverseType) {
					case 1:
					System.out.print("\nPreorder Traversal: ");
					preOrder(root);
					break;
					case 2:
					System.out.print("\nInorder Traversal: ");
					inOrder(root);
					break;
					case 3:
					System.out.println("\nPostorder Traversal: ");
					postOrder(root);
					break;
				}
				System.out.println();
			}
			
			private void preOrder(Node localRoot) {
				if (localRoot != null) {
					System.out.print(localRoot.iData + " ");
					preOrder(localRoot.leftChild);
					preOrder(localRoot.rightChild);
				}
			}
			
			private void inOrder(Node localRoot) {
				if (localRoot != null) {
					inOrder(localRoot.leftChild);
					System.out.print(localRoot.iData + " ");
					inOrder(localRoot.rightChild);
				}
			}
			
			private void postOrder(Node localRoot) {
				if (localRoot != null) {
					postOrder(localRoot.leftChild);
					postOrder(localRoot.rightChild);
					System.out.print(localRoot.iData + " ");
				}
			}
			
			public void displayTree() {
				Stack<Node> globalStack = new Stack<Node>();
				globalStack.push(root);
				int nBlanks = 32;
				boolean isRowEmpty = false;
				System.out.println("--------------------------------------------------");
				while (isRowEmpty == false) {
					Stack<Node> localStack = new Stack<Node>();
					isRowEmpty = true;
					
					for (int j = 0; j < nBlanks; j++) {
						System.out.print(" ");
					}
					
					while (globalStack.isEmpty() == false) {
						Node temp = globalStack.pop();
						if (temp != null) {
							System.out.print(temp.iData);
							localStack.push(temp.leftChild);
							localStack.push(temp.rightChild);
							
							if (temp.leftChild != null || temp.rightChild != null) {
								isRowEmpty = false;
							}
						} else {
							System.out.print("  ");
							localStack.push(null);
							localStack.push(null);
						}
						
						for (int j = 0; j < nBlanks * 2 - 2; j++) {
							System.out.print(" ");
						}
					} // end while globalStack not empty
					System.out.println();
					nBlanks /= 2;
					while (localStack.isEmpty() == false) {
						globalStack.push(localStack.pop());
					}
				} // end while isRowEmpty is false
				System.out.println("---------------------------------------------------");
			} // end displayTree()
			
			public Node minimum() {
				Node current;
				Node last = null;
				current = root;
				while (current != null) {
					last = current;
					current = current.leftChild;
				}
				return last;
			}
			
			public Node maximum() {
				Node current;
				Node last = null;
				current = root;
				while (current != null) {
					last = current;
					current = current.rightChild;
				}
				return last;
			}
		} // end class Tree
	\end{lstlisting}
	
	\textbf{TreeApp.java}
	\begin{lstlisting}
		package p1;
		
		import java.io.BufferedReader;
		import java.io.IOException;
		import java.io.InputStreamReader;
		
		
		public class TreeApp {
			
			public static void main(String[] args) throws IOException {
				int value;
				Tree theTree = new Tree();
				
				theTree.insert(50, 1.5);
				theTree.insert(25, 1.2);
				theTree.insert(75, 1.7);
				theTree.insert(12, 1.5);
				theTree.insert(37, 1.2);
				theTree.insert(43, 1.7);
				theTree.insert(30, 1.7);
				theTree.insert(33, 1.5);
				theTree.insert(87, 1.7);
				theTree.insert(93, 1.5);
				theTree.insert(97, 1.5);
				
				while (true) {
					System.out.print("Enter first letter of show, ");
					System.out.print("insert, find, delete, or traverse: ");
					int choice = getChar();
					switch (choice) {
						case 's':
						theTree.displayTree();
						break;
						case 'i':
						System.out.print("Enter value to insert: ");
						value = getInt();
						theTree.insert(value, value + 0.9);
						break;
						case 'f':
						System.out.print("Enter value to find: ");
						value = getInt();
						Node found = theTree.find(value);
						if(found != null) {
							System.out.print("found: ");
							found.displayNode();
							System.out.println("\n");
						} else {
							System.out.print("Could not find ");
							System.out.print(value + '\n');
						}
						break;
						case 'd':
						System.out.print("Enter value to delete: ");
						value = getInt();
						boolean didDelete = theTree.delete(value);
						if(didDelete) {
							System.out.print("Deleted " + value + "\n");
						} else {
							System.out.print("Could not delete ");
							System.out.print(value + '\n');
						}
						break;
						case 't':
						System.out.print("Enter type 1, 2 or 3: ");
						value = getInt();
						theTree.traverse(value);
						break;
						default:
						System.out.print("Invalid entry.\n");
					} // end switch
				}    // end while
			}    // end main()
			
			public static String getString() throws IOException{
				InputStreamReader isr = new InputStreamReader(System.in);
				BufferedReader br = new BufferedReader(isr);
				String s = br.readLine();
				return s;
			}
			
			public static char getChar() throws IOException {
				String s = getString();
				return s.charAt(0);
			}
			
			public static int getInt() throws IOException {
				String s = getString();
				return Integer.parseInt(s);
			}
		}
	\end{lstlisting}
	
\end{document}